\chapter{Database Connections and Transactions}
\label{ch:db-connection}

Kiss provides a database abstraction layer that supports five database engines through a unified API.  This chapter covers the \texttt{Connection}\index{Connection}, \texttt{Command}\index{Command}, and \texttt{Cursor}\index{Cursor} classes that form the foundation of all database access.

\section{The Connection Class}
\label{sec:db-connection}

\texttt{Connection} wraps a JDBC connection with convenience methods for queries, transactions, and metadata access.  In a service method, the framework provides a \texttt{Connection} as the \texttt{db} parameter --- you rarely need to create one yourself.

\begin{lstlisting}[style=groovy]
void myMethod(JSONObject injson, JSONObject outjson,
              Connection db, ProcessServlet servlet) {
    // db is ready to use --- connection pooling, auto-commit
    // off, and transaction management are all handled by
    // the framework
    Record emp = db.fetchOne(
        "select * from employee where employee_id = ?", empId)
}
\end{lstlisting}

\subsection{Creating Standalone Connections}

For batch jobs, utilities, or code outside the service framework, you can create a connection directly:

\begin{lstlisting}[style=java]
try (Connection db = new Connection(
        ConnectionType.PostgreSQL,
        "localhost", null, "mydb", "user", "password")) {
    // use db ...
    db.commit();
}  // connection auto-closes via try-with-resources
\end{lstlisting}

\noindent
Available \texttt{ConnectionType}\index{ConnectionType} values: \texttt{PostgreSQL}, \texttt{MySQL}, \texttt{SQLite}, \texttt{MicrosoftServer}, \texttt{Oracle}.

\section{Supported Databases}
\label{sec:db-databases}

Kiss supports five database engines.  The database type is specified in \texttt{application.ini} and affects SQL generation for features like \texttt{LIMIT} and pagination:

\begin{longtable}{lll}
\toprule
\textbf{Database} & \textbf{ConnectionType} & \textbf{LIMIT Syntax} \\
\midrule
\endhead
PostgreSQL & \texttt{PostgreSQL} & \texttt{LIMIT n} \\
MySQL & \texttt{MySQL} & \texttt{LIMIT n} \\
SQLite & \texttt{SQLite} & \texttt{LIMIT n} \\
SQL Server & \texttt{MicrosoftServer} & \texttt{TOP n} \\
Oracle & \texttt{Oracle} & \texttt{OFFSET ... FETCH} \\
\bottomrule
\end{longtable}

\section{Executing SQL}
\label{sec:db-execute}

The \texttt{execute()} method runs INSERT, UPDATE, and DELETE statements:

\begin{lstlisting}[style=groovy]
// Update
db.execute(
    "UPDATE employee SET active = ? WHERE employee_id = ?",
    "N", empId)

// Insert
db.execute("INSERT INTO log (message, created) VALUES (?, ?)",
           msg, new Date())

// Delete
db.execute("DELETE FROM temp_data WHERE created < ?",
    cutoffDate)
\end{lstlisting}

\begin{warnbox}
Always use \texttt{?} parameter placeholders instead of string concatenation.  This prevents SQL injection and handles type conversion automatically.
\end{warnbox}

\subsection{Existence Checks}

\begin{lstlisting}[style=groovy]
if (db.exists("SELECT * FROM employee WHERE email = ?", email))
    throw new UserException("Email already in use")
\end{lstlisting}

\subsection{Row Counting}

\begin{lstlisting}[style=groovy]
long total = db.fetchCount(
    "SELECT * FROM employee WHERE active = 'Y'")
\end{lstlisting}

\noindent
\texttt{fetchCount()} wraps the query in \texttt{SELECT COUNT(*) FROM (...)} --- use it primarily for pagination display, as it executes the full query internally.

\section{Fetching Data}
\label{sec:db-fetch}

Kiss provides several methods for retrieving data, each suited to a different use case.

\subsection{Single Record}

\begin{lstlisting}[style=groovy]
Record emp = db.fetchOne(
    "SELECT * FROM employee WHERE employee_id = ?", id)
if (emp == null)
    throw new UserException("Employee not found")
String name = emp.getString("first_name")
\end{lstlisting}

\noindent
\texttt{fetchOne()} returns the first record or \texttt{null} if no rows match.  The returned \texttt{Record} can be updated or deleted (see Chapter~\ref{ch:db-crud}).

\subsection{All Records}

\begin{lstlisting}[style=groovy]
List<Record> emps = db.fetchAll(
    "SELECT * FROM employee WHERE active = 'Y'")
for (Record emp : emps) {
    String name = emp.getString("last_name")
    // process each record
}
\end{lstlisting}

\noindent
Returns an empty list (not \texttt{null}) if no rows match.

\subsection{Limited Records}

\begin{lstlisting}[style=groovy]
// Fetch at most 100 records
List<Record> page = db.fetchAll(100,
    "SELECT * FROM employee ORDER BY last_name")
\end{lstlisting}

\subsection{Paginated Records}

\begin{lstlisting}[style=groovy]
// Fetch page 2 (zero-based), 25 records per page
List<Record> page = db.fetchAll(2, 25,
    "SELECT * FROM employee ORDER BY last_name")
\end{lstlisting}

\subsection{Fetching as JSON}

For the common case of returning query results directly to the frontend:

\begin{lstlisting}[style=groovy]
// All records as JSON array
outjson.put("employees", db.fetchAllJSON(
    "SELECT employee_id, first_name, last_name FROM employee"))

// Single record as JSON object
outjson.put("employee", db.fetchOneJSON(
    "SELECT * FROM employee WHERE employee_id = ?", id))
\end{lstlisting}

\noindent
\texttt{fetchAllJSON()} also supports \texttt{max} and \texttt{page/max} variants.

\begin{tipbox}
Use \texttt{fetchAllJSON()} and \texttt{fetchOneJSON()} when the database column names match what the frontend expects.  This eliminates the need to manually build JSON objects from Records.
\end{tipbox}

\section{Cursors for Large Result Sets}
\label{sec:db-cursor}

When a result set is too large to hold in memory, use a \texttt{Cursor} to iterate row by row:

\begin{lstlisting}[style=groovy]
try (Cursor cursor = db.query(
        "SELECT * FROM large_table ORDER BY id")) {
    while (cursor.isNext()) {
        Record r = cursor.getRecord()
        String name = r.getString("name")
        // process one row at a time
    }
}
\end{lstlisting}

\noindent
Despite the name, Kiss cursors are \emph{not} database cursors.  The framework downloads the entire result set into a local cache (in memory for small sets, on disk for large ones) and then closes the database cursor.  This means:

\begin{itemize}
  \item Multiple cursors can be open simultaneously without database performance issues.
  \item Nested cursor operations work correctly.
  \item Memory usage is controlled by the caching strategy, not the result set size.
\end{itemize}

\noindent
The \texttt{query()} method also supports \texttt{max} and \texttt{page/max} parameters:

\begin{lstlisting}[style=groovy]
// At most 1000 rows
Cursor cursor = db.query(1000, "SELECT * FROM audit_log")

// Page 3, 50 rows per page
Cursor cursor = db.query(3, 50, "SELECT * FROM audit_log")
\end{lstlisting}

\section{The Command Class}
\label{sec:db-command}

\texttt{Command} provides fine-grained control over prepared statements.  Use it when you need to execute the same SQL repeatedly or need a separate statement context:

\begin{lstlisting}[style=groovy]
try (Command cmd = db.newCommand()) {
    cmd.execute("INSERT INTO log (message) VALUES (?)", msg)
}
\end{lstlisting}

\noindent
\texttt{Command} supports all the same methods as \texttt{Connection} --- \texttt{execute()}, \texttt{fetchOne()}, \texttt{fetchAll()}, \texttt{fetchAllJSON()}, \texttt{query()}, \texttt{exists()} --- but against its own prepared statement.  This is useful when you need multiple active statements simultaneously.

\begin{lstlisting}[style=groovy]
try (Command selectCmd = db.newCommand()
     Command updateCmd = db.newCommand()) {
    Cursor cursor = selectCmd.query(
        "SELECT * FROM employee WHERE needs_update = 'Y'")
    while (cursor.isNext()) {
        Record r = cursor.getRecord()
        updateCmd.execute(
            "UPDATE employee SET processed = 'Y' " +
            "WHERE employee_id = ?",
            r.getInt("employee_id"))
    }
    db.commit()
}
\end{lstlisting}

\noindent
\texttt{Command} also provides \texttt{newQueryBuilder()} for using the \texttt{QueryBuilder} (Chapter~\ref{ch:query-builder}) with a specific command instance.  This is essential when iterating over one query's results while executing another query on the same connection:

\begin{lstlisting}[style=groovy]
try (Command cmd1 = db.newCommand()
     Command cmd2 = db.newCommand()) {
    Cursor cursor = cmd1.newQueryBuilder()
        .select("employee.first_name", "employee.last_name")
        .where("employee.active = ?", "Y")
        .query()
    while (cursor.isNext()) {
        Record emp = cursor.getRecord()
        List<Record> tasks = cmd2.newQueryBuilder()
            .select("task.description", "task.due_date")
            .where("task.employee_id = ?",
                   emp.getString("employee_id"))
            .fetchAll()
        // process tasks...
    }
}
\end{lstlisting}

\section{Transactions}
\label{sec:db-transactions}

Kiss connections always have auto-commit disabled.  This means:

\begin{itemize}
  \item Changes are not visible to other connections until \texttt{commit()} is called.
  \item The framework calls \texttt{commit()} automatically when a service method returns successfully.
  \item The framework calls \texttt{rollback()} automatically when a service method throws an exception.
  \item After each \texttt{commit()} or \texttt{rollback()}, a new transaction begins automatically.
\end{itemize}

\begin{lstlisting}[style=groovy]
// Usually you don't need to call commit() explicitly. The
// framework does it for you at the end of each service call.

// Exception: long-running batch operations
int count = 0
for (Record rec : records) {
    db.execute("UPDATE ...", rec.getInt("id"))
    if (++count % 100 == 0)
        db.commit()  // commit every 100 rows
}
db.commit()  // commit remaining rows
\end{lstlisting}

\begin{warnbox}
If you call \texttt{db.commit()} partway through a service method and then throw an exception, the committed changes cannot be rolled back.  Use mid-method commits only for batch operations where partial progress is acceptable.
\end{warnbox}

\section{Parameterized Queries}
\label{sec:db-params}

All query and execute methods accept \texttt{?} parameter placeholders:

\begin{lstlisting}[style=groovy]
db.fetchAll(
    "SELECT * FROM employee WHERE dept = ? AND active = ?",
    deptId, "Y")
\end{lstlisting}

\noindent
Parameters are bound using JDBC prepared statements.  This:

\begin{itemize}
  \item Prevents SQL injection attacks.
  \item Handles type conversion automatically (strings, numbers, dates).
  \item Allows the database to cache query plans.
\end{itemize}

\noindent
Date parameters are handled automatically.  You can pass \texttt{java.util.Date}, \texttt{java.util.Calendar}, \texttt{ZonedDateTime}, \texttt{LocalDateTime}, \texttt{LocalDate}, or \texttt{LocalTime} objects, and Kiss will convert them to the appropriate JDBC type.

\section{Database-Specific SQL}
\label{sec:db-specific}

When you need database-portable row limiting, pass a maximum row count to the fetch methods instead of writing \texttt{LIMIT} or \texttt{TOP} directly.  The \texttt{max}-argument overloads (and \texttt{QueryBuilder.limit(int)}) generate the correct syntax for the connected database:

\begin{lstlisting}[style=groovy]
// Limit to max rows (database-portable)
List<Record> recs = db.fetchAll(100,
    "SELECT * FROM employee ORDER BY last_name")
\end{lstlisting}

\noindent
Internally Kiss produces the correct syntax for the connected database:
\begin{itemize}
  \item PostgreSQL/MySQL/SQLite: appends \texttt{LIMIT 100}
  \item SQL Server: inserts \texttt{TOP 100} after \texttt{SELECT}
  \item Oracle: appends \texttt{OFFSET 0 ROWS FETCH NEXT 100 ROWS ONLY}
\end{itemize}

\section{Connection Metadata}
\label{sec:db-metadata}

\texttt{Connection} provides several metadata methods:

\begin{lstlisting}[style=groovy]
// Check if a table exists
if (db.tableExists("employee")) { ... }

// Get the width of a VARCHAR column (cached, exact)
int maxLen = db.getColumnWidth("employee", "last_name")

// Get primary key column names
List<String> pkCols = db.getPrimaryColumns("employee")

// Get single primary key column name
String pkCol = db.getPrimaryColumnName("employee")

// Get the database type
ConnectionType type = db.getDBType()
\end{lstlisting}

\paragraph{Column widths.} \texttt{getColumnWidth(table, column)}\index{Connection!getColumnWidth} returns the declared width of a column in characters, read from the database metadata, so a value can be checked against the schema instead of against a constant repeated in Java: \texttt{value.length() > db.getColumnWidth("person", "fname")}.  Each answer is read once and then remembered for the life of the JVM, per schema (or catalog), table and column; the cache is thread-safe and shared by all connections.  Names are case-insensitive, and the match is exact --- a \texttt{\_} in a table or column name is an ordinary character, not the single-character wildcard that \texttt{DatabaseMetaData} patterns use.  A column with no limit (text, or a character type declared without a length) and a column that does not exist both answer \texttt{Connection.UNLIMITED\_COLUMN\_WIDTH} (\texttt{Integer.MAX\_VALUE}), so the comparison above never rejects anything for them; a missing column never throws and is not remembered, so a column added later is picked up.  An \texttt{SQLException} raised by the database while reading the metadata is not hidden: it propagates, because treating an outage as ``unlimited'' would silently switch the check off.  \texttt{getColumnSize()} is the older, uncached call: it treats the names as patterns and answers $-1$ for a missing column.

\paragraph{Metadata is scoped to the current schema.} Every metadata read --- column definitions, column sizes, and primary keys --- is scoped to the connection's current schema (the current catalog on MySQL). In a multi-schema database, such as a schema-per-tenant layout, a table of the same name in another schema never influences the result. Anything derived from metadata follows the same rule; for example, the inline threshold of a hybrid column (Section~\ref{sec:file-external-hybrid}) always reflects the current schema's column definition.

\section{Connection Pooling}
\label{sec:db-pooling}

Kiss uses C3P0\index{C3P0} for connection pooling.  Pool parameters are configured automatically based on CPU core count:

\begin{itemize}
  \item \textbf{Minimum pool size}: max(2, number of cores)
  \item \textbf{Initial pool size}: max(minimum, cores $\times$ 2)
  \item \textbf{Maximum pool size}: cores $\times$ 4, minimum 20
  \item \textbf{Idle timeout}: 30 minutes
  \item \textbf{Test period}: Every 5 minutes on idle connections
  \item \textbf{Checkout timeout}: 10 seconds
\end{itemize}

\noindent
These defaults work well for most applications.  For high-load production deployments, see Chapter~\ref{ch:deployment} for tuning guidance.

\section{Per-Request Connection Preparation}
\label{sec:db-request-preparer}

Every REST request that carries a database connection gets that connection handed to the service method fully ready to use.  Sometimes ``ready to use'' involves per-request state --- the canonical example is a \textbf{multi-tenant} application that stores each tenant in its own database schema and must point every request's connection at the right one (for example, the PostgreSQL \texttt{search\_path}) before any service code runs.

\texttt{RequestConnectionPreparer}\index{RequestConnectionPreparer} is the framework hook for exactly this.  An application registers one implementation, typically from \texttt{KissInit.groovy}:

\begin{lstlisting}[style=groovy]
MainServlet.setRequestConnectionPreparer(
    new MyConnectionPreparer())
\end{lstlisting}

\noindent
The interface has two methods:

\begin{lstlisting}[style=java]
public interface RequestConnectionPreparer {
    /** Called after authentication, before the web service
     *  method runs. */
    void prepare(Connection db,
        UserData userData) throws Exception;

    /** Called when a request connection is closed, before the
     *  underlying connection returns to the pool.  Default:
     *  no action. */
    default void release(Connection db) {
    }
}
\end{lstlisting}

The contract:

\begin{itemize}
  \item \texttt{prepare} runs for every REST request that carries a database connection, \emph{after} authentication has completed and \emph{before} the service method executes.  \texttt{userData} is the authenticated user --- or \texttt{null} when the method is allowed without authentication (Chapter~\ref{ch:authentication}) and no valid session accompanies the request.
  \item \textbf{Fail-closed:} if \texttt{prepare} throws, the request is aborted with an error response and the service method never runs.  An authenticated session that cannot be mapped to a valid tenant therefore reaches no business data.
  \item \texttt{release} runs when the request's connection is closed, before the underlying pooled connection is returned --- use it to clear whatever \texttt{prepare} set (reset the search path, clear thread-local state), so pooled connections never leak state from one consumer to the next.
  \item Connections obtained \emph{outside} the request path --- cron tasks, \texttt{MainServlet.openNewConnection()} --- are \textbf{not} passed to \texttt{prepare}; such consumers must establish any needed connection state themselves.
\end{itemize}

A multi-tenant implementation, condensed:

\begin{lstlisting}[style=java]
public class MyConnectionPreparer
        implements RequestConnectionPreparer {
    @Override
    public void prepare(Connection db,
        UserData userData) throws Exception {
        if (userData == null) {
            //  Allowed-without-authentication: public
            //  schema only.
            setSearchPath(db, "public");
            return;
        }
        // set at login
        String tenantId = (String) userData.getUserData(
            "tenantId");
        if (tenantId == null)
            throw new Exception(
                "No tenant is associated with this session");
        // every service is now tenant-neutral
        setSearchPath(db, tenantId);
    }

    @Override
    public void release(Connection db) {
        // never return a tenant-pointed
        setSearchPath(db, "public");
    }
}
\end{lstlisting}

\noindent
The payoff of this design is that \emph{no web service ever mentions a tenant}: services read and write tables normally, and the schema selection happens once, in one place, for every request.  When no preparer is registered the hook costs nothing.
